\documentclass{article} % This command is used to set the type of
% document you are working on such as an article, book, or presenation
\usepackage{
  amsmath,      % Math Environments
  amssymb,      % Extended Symbols
  enumerate,        % Enumerate Environments
  graphicx,      % Include Images
  lastpage,      % Reference Lastpage
  multicol,      % Use Multi-columns
  multirow      % Use Multi-rows
}

\usepackage{geometry} % This package allows the editing of the page layout
\usepackage{amsmath}  % This package allows the use of a large range
% of mathematical formula, commands, and symbols
\usepackage{graphicx}  % This package allows the importing of images

\newcommand{\question}[2][]{
  \begin{flushleft}
    \textbf{Question #1}: \textit{#2}

\end{flushleft}}
\newcommand{\sol}{\textbf{Solution}:} %Use if you want a boldface solution line
\newcommand{\maketitletwo}[2][]{
  \begin{center}
    \Large{\textbf{Assignment #1}

    MATH 4362} % Name of course here
    \vspace{5pt}

    \normalsize{N. Ohayon Rozanes  % Your name here

    \today}        % Change to due date if preferred
    \vspace{15pt}

\end{center}}
\begin{document}
\maketitletwo[5]  % Optional argument is assignment number
%Keep a blank space between maketitletwo and \question[1]

\question[1]{}
\[
  \begin{cases}
    u_{xx} + u_{yy} = 0 \\
    u(x,0) = \sin^3x\\
    u(x,\pi) = u(0, y) = u(\pi, y) = 0
  \end{cases}
\]

Assuming a seperable ansatz:

$$
u(x,y) = v(x)w(y)
$$

It follows from substituting into the original equation that:

\[
  \frac{v''}{v} = -\frac{w''}{w} = \lambda
\]

Which yields the two ODEs:

\[
  v'' - v\lambda = 0
\]
and
\[
  w'' + w\lambda = 0
\]

$v$'s ODE is a second order differential equation with imaginary
roots, and has the solution:

$$
v(x) = c_1 \cos(\sqrt{-\lambda}x) + c_2\sin(\sqrt{-\lambda}x)
$$

Conditioned on $u(0,y) = 0$ we have that:

\[
  v(0) = 0 \implies c_1\cos(0) = 0 \implies c_1 = 0
\]

similarily, after conditioning on $u(\pi, y) = 0$, we have that
\[
  v(\pi) = 0 \implies c_2\sin(\sqrt{-\lambda}\pi) = 0
\]

Which implies either $c_2 = 0$, which leads to trivial solutions, or
that $\sqrt{-\lambda} \in \mathbb{N}$, that is $\lambda = -n^2$ for
any $n \in \mathbb{N}$.

$w$'s ODE is a second order differential equation only with real
roots, namely: $\sqrt{-\lambda}$ and $-\sqrt{-\lambda}$ and thus has
the solution:

\[
  w = c_1 e^{\sqrt{-\lambda}y} + c_2e^{-\sqrt{-\lambda}y}
\]

Finally conditioning on $u(x,\pi) = 0$ we have that
\[
  w(\pi) = c_1 e^{\sqrt{-\lambda}\pi} + c_2e^{-\sqrt{-\lambda}\pi} = 0
\]

which after some calculations

\[
  w(y) = c_n \sinh(\sqrt{-\lambda}(\pi-y)) = c_n\sinh(\sqrt{n(\pi-y)})
\]

Then the function $u$ is, for any $n$

$$
u = c_n \sinh(n(\pi-y))\sin(nx)
$$

Since Laplace is linear, we can some any number of such solutions to
obtain a solution:
$$
u = \sum c_n \sinh(n(\pi-y))\sin(nx)
$$

Finally conditioning on the boundary at $u(x,0) = sin^3x$:

\[
  u(x,0) = \sin^3x = \sum c_n \sinh(n\pi)\sin(nx)
\]

Setting $d_n = c_n \sinh(n\pi)$, we use the fourier series to find
$d_n$, given by the integral:

\[
  d_n = \frac{2}{\pi}\int_0^\pi \sin^3x\sin(nx)\,dx
\]

To evaluate the above, first consider the power reduction formula
\[
  \sin^3x = \frac{3\sin x - \sin 3x}{4}
\]

then the integrand becomes

$$
\sin^3x\sin(nx) = \frac{3\sin(x)\sin(nx) - \sin(3x)\sin(nx)}{4}
$$

Since trigonometric functions are orthogonal in $L[0,\pi]$, either
term of the above disappears if the frequencies do not match, that
is, the only $n$ of interest are $n = 1$ and $n = 3$. With that in
mind, we solve the following integrals:

\[
  n = 1:\quad \int_0^\pi \sin^3x\sin(x)\,dx
  = \frac{1}{4}\left(3\int_0^\pi \sin^2x\,dx
  - \int_0^\pi \sin(3x)\sin(x)\,dx\right)
  = \frac{1}{4}\left(3\cdot\frac{\pi}{2}-0\right)
  = \frac{3\pi}{8}
\]

\[
  n = 3:\quad \int_0^\pi \sin^3x\sin(3x)\,dx
  = \frac{1}{4}\left(3\int_0^\pi \sin(x)\sin(3x)\,dx
  - \int_0^\pi \sin^2(3x)\,dx\right)
  = \frac{1}{4}\left(0-\frac{\pi}{2}\right)
  = -\frac{\pi}{8}
\]

and for $n \ne 1,3$ the integral is $0$. Therefore
\[
  d_n = \frac{2}{\pi}\int_0^\pi \sin^3x\sin(nx)\,dx
  =
  \begin{cases}
    \frac{3}{4}, & n = 1,\\
    -\frac{1}{4}, & n = 3,\\
    0, & n \ne 1,3.
  \end{cases}
\]

Recalling $d_n = c_n \sinh(n\pi)$, we obtain
\[
  c_1 = \frac{3}{4\sinh(\pi)},\quad
  c_3 = -\frac{1}{4\sinh(3\pi)},\quad
  c_n = 0\ \text{for}\ n \ne 1,3.
\]
Thus the final solution is
\[
  u(x,y)
  = \frac{3}{4\sinh(\pi)}\sinh(\pi-y)\sin x
  - \frac{1}{4\sinh(3\pi)}\sinh(3(\pi-y))\sin(3x).
\]

\question[2]{}

The solution to this problem would be the sum of the solution of the
two Dirichlet problems:
$$
\begin{cases}
  u_{xx} + u_{yy} = 0 \\
  u(x,0) = \sin\pi x\\
  u(x,\pi) = u(0, y) = u(\pi, y) = 0
\end{cases}
\begin{cases}
  u_{xx} + u_{yy} = 0 \\
  u(0,y) = \sin\pi y\\
  u(x,\pi) = u(x, 0) = u(\pi, y) = 0
\end{cases}
$$

By an almost identical analysis to question $1$, in which we utilize
the seperable anstaz and condition on the boundary functions, the above two
Dirichlet problems have solutions of the form

$$
\sum c_n\sinh(n\pi(1-y))\sin(n\pi x)
$$
and
$$
\sum c_n\sinh(n\pi(1-x))\sin(n\pi y)
$$

respectively. Focusing on just the first solution and conditioning on
$u(x,0) = \sin \pi x$

$$
\sin \pi x = \sum c_n\sinh(n\pi(1-y))\sin(n\pi x)
$$

which gives
\[
  c_1 = \frac{1}{\sinh(\pi)},\quad c_n = 0\ \text{for}\ n \ne 1.
\]
Thus the first solution is
\[
  u_1(x,y) = \frac{1}{\sinh(\pi)}\sinh(\pi(1-y))\sin(\pi x).
\]

For the second problem, the same calculation yields
\[
  u_2(x,y) = \frac{1}{\sinh(\pi)}\sinh(\pi(1-x))\sin(\pi y).
\]

Therefore the final solution is
\[
  u(x,y) = \frac{1}{\sinh(\pi)}\left[\sinh(\pi(1-y))\sin(\pi x)
  + \sinh(\pi(1-x))\sin(\pi y)\right].
\]
\question[3]{}
\[
  \begin{cases}
    u_{xx} + u_{yy} = 0 \\
    u = x^3 \quad x^2 + y^2 = 1
  \end{cases}
\]

Consider that the general solution for the Laplace equation on a unit disk
in polar coordinates is:

\[
  U = \frac{a_0}{2} + \sum a_n r^n \cos(n\theta) + b_n r^n\sin(n\theta)
\]

The boundary condition in polar coordinates becomes

\[
  U(1, \theta) = H(\theta)  = \cos^3\theta
\]

Since each $\cos$ is in itself a $2\pi$ periodic function, $\cos^3$
is as well. In calculating the fourier coefficients for the solution notice that
the integrand of $b_n$ is odd:
$$
b_n = \frac{1}{\pi}\int_{-\pi}^\pi \cos^3\theta\sin\theta d\theta
$$

Thus, it disappears over a symmetric domain, and therefore $b_n = 0$
for all $n$. A similar trick to question $1$ is used to compute the
values of $a_n$, where first the power reducing identity is used
then, using orthogonality of trigonmetric functions in $L[-\pi,\pi]$,
we discard all non matching frequencies

\[
  \begin{aligned}
    \cos^3\theta \cos(n\theta)
    &= \cos\theta\left(\frac{1+\cos(2\theta)}{2}\right)\cos(n\theta) \\
    &= \frac{\cos\theta+\cos\theta\cos(2\theta)}{2}\cos(n\theta) \\
    &=
    \left(\frac{\cos\theta}{2}+\frac{\cos(3\theta)+\cos(-\theta)}{4}\right)\cos(n\theta)
    \\
    &=
    \left(\frac{3\cos\theta}{4}+\frac{\cos(3\theta)}{4}\right)\cos(n\theta) \\
    &= \frac{3\cos\theta\cos(n\theta)+\cos(3\theta)\cos(n\theta)}{4}
  \end{aligned}
\]

Using orthogonality on $[-\pi,\pi]$, only $n=1$ and $n=3$ contribute.
Thus
\[
  a_n = \frac{1}{\pi}\int_{-\pi}^{\pi}\cos^3\theta\cos(n\theta)\,d\theta
  =
  \begin{cases}
    \frac{3}{4}, & n = 1,\\
    \frac{1}{4}, & n = 3,\\
    0, & n \ne 1,3,
  \end{cases}
\]
and $a_0 = 0$ while $b_n = 0$ for all $n$. Therefore the harmonic
solution in the disk is
\[
  U(r,\theta) = \frac{3}{4}r\cos\theta + \frac{1}{4}r^3\cos(3\theta).
\]
\question[4]{}

\[
  \begin{cases}
    u_{xx} + u_{yy} = 0 \\
    u = x^3 \quad x^2 + y^2 = 4
  \end{cases}
\]
A similar process to question $3$ follows, the above becomes
$$
\begin{cases}
  u_{rr} + u_{\theta\theta} = 0 \\
  U = H(\theta) \quad r = 4
\end{cases}
$$

Which has the following general form:

\[
  U = \frac{a_0}{2} + \sum a_n r^n \cos(n\theta) + b_n r^n\sin(n\theta)
\]

The boundary condition, in polar coordinates becomes

\[
  U(4, \theta) = H(\theta) = 4^4\cos^4\theta = 256 \cos^4\theta
\]

Then we have that

\[
  U(4, \theta) = \frac{a_0}{2} + \sum a_n 4^n \cos(n\theta) + b_n
  4^n\sin(n\theta) = 256\cos^4\theta
\]
Setting $c_n = 4^na_n$ and $d_n = 4^nb_n$,
\[
  \frac{c_0}{2} + \sum c_n\cos(n\theta) + d_n\sin(n\theta)
  = 256\cos^4\theta
\]
Is just the typical Fourier series set up and clearly $d_n = 0$. Then,
working the boundary
condition to an integrable form:
\[
  \begin{aligned}
    \cos^4\theta
    &= \cos^2\theta\cos^2\theta \\
    &= \left(\frac{1+\cos 2\theta}{2}\right)\left(\frac{1+\cos
    2\theta}{2}\right) \\
    &= \frac{\cos^2 2\theta + 2\cos 2\theta + 1}{4} \\
    &= \frac{1+\cos 4\theta}{8} + \frac{\cos 2\theta}{2} + \frac{1}{4} \\
    &= \frac{\cos 4\theta}{8} + \frac{\cos 2\theta}{2} + \frac{3}{8}
  \end{aligned}
\]
Multiplying through by $256\cos n\theta$
\[
  \begin{aligned}
    256(\cos^4\theta)\cos(n\theta)
    &= 32\cos 4\theta\cos(n\theta)
    + 128\cos 2\theta\cos(n\theta)
    + 96\cos(n\theta) \\
  \end{aligned}
\]

Since $H(\theta)$ is even, $d_n = 0$. Using orthogonality on
$[-\pi,\pi]$,
\[
  c_n = \frac{1}{\pi}\int_{-\pi}^{\pi}256\cos^4\theta\cos(n\theta)\,d\theta
  =
  \begin{cases}
    192, & n = 0,\\
    8, & n = 2,\\
    \frac{1}{8}, & n = 4,\\
    0, & n \ne 0,2,4.
  \end{cases}
\]

Recall that $a_n = \frac{c_n}{4^n}$, and the harmonic solution in the disk is
\[
  U(r,\theta) = 96 + 8r^2\cos(2\theta) + \frac{1}{8}r^4\cos(4\theta).
\]

\end{document}
